\subsection{Quasi-invariant measures on a homogeneous space}

Lemma 3. --- Let G be a locally compact group, \(\mu\) a left Haar measure on G, \(\nu\) and \(\nu'\) two nonzero quasi-invariant measures on G. If, for every \(s \in G\) , the densities of \(\boldsymbol{\gamma}(s)\nu\) with respect to \(\nu\) and of \(\boldsymbol{\gamma}(s)\nu'\) with respect to \(\nu'\) are equal locally \(\mu\) -almost everywhere, then \(\nu\) and \(\nu'\) are proportional.

Write \(\nu = \rho \cdot \mu\) , \(\nu' = \rho' \cdot \mu\) , where \(\rho, \rho'\) are locally \(\mu\) -integrable functions on \(G\) and are everywhere nonzero (§1, No.~9, Prop. 11). For every \(s \in G\) ,

\[
\boldsymbol {\gamma} (s) \nu = (\boldsymbol {\gamma} (s) \rho) \cdot \mu , \quad \boldsymbol {\gamma} (s) \nu^ {\prime} = (\boldsymbol {\gamma} (s) \rho^ {\prime}) \cdot \mu ,
\]

and the hypothesis implies that \(\rho^{-1} \cdot \pmb{\gamma}(s) \rho = \rho'^{-1} \cdot \pmb{\gamma}(s) \rho'\) locally \(\mu\) -almost everywhere. Set \(\sigma = \rho'/\rho\) , which is a \(\mu\) -measurable function on \(G\) . For every \(s \in G\) , \(\pmb{\gamma}(s) \sigma = \sigma\) locally \(\mu\) -almost everywhere. Therefore \(\sigma\) is equal to a constant locally \(\mu\) -almost everywhere, by Cor. 2 of Prop. 6 applied with \(X = H = G\) .

Let G be a locally compact group, H a closed subgroup of G. Consider the homogeneous space G/H of left cosets with respect to H, on which G operates continuously on the left. We are going to show that there exists one and only one class of nonzero quasi-invariant measures on G/H.

Note that H operates on G continuously and properly by right translations; and the quotient space, which is none other than G/H, is paracompact (GT, III, §4, No.~6, Prop. 13). We may therefore apply the results of Nos. 1 to 4, with X = G. We thus have mappings \(f \mapsto f^{b}\) of \(\mathcal{H}(G)\) onto \(\mathcal{K}(G/H)\) , and \(\lambda \mapsto \lambda^{\sharp}\) of \(\mathcal{M}(G/H)\) into \(\mathcal{M}(G)\) (once a left Haar measure \(\beta\) on H has been fixed). The fact that G operates on the left in G/H gives rise to a supplementary property:

\[
\boldsymbol {\gamma} _ {\mathrm{G/H}} (s) \cdot f ^ {b} = \left(\boldsymbol {\gamma} _ {\mathrm{G}} (s) \cdot f\right) ^ {b} \quad \left(s \in \mathrm{G}, f \in \mathscr {K} (\mathrm{G})\right)\tag{15}
\]

\[
\left(\boldsymbol {\gamma} _ {\mathrm{G/H}} (s) \cdot \lambda\right) ^ {\sharp} = \boldsymbol {\gamma} _ {\mathrm{G}} (s) \cdot \lambda^ {\sharp} \quad \left(s \in \mathrm{G}, \lambda \in \mathscr {M} (\mathrm{G/H})\right).\tag{16}
\]

Indeed, for every \(x \in G\) ,

\[
\begin{array}{r l} & {\left(\boldsymbol {\gamma} _ {\mathrm{G/H}} (s) \cdot f ^ {b}\right) (\pi (x)) = f ^ {b} \big (s ^ {- 1} \pi (x) \big) = f ^ {b} \big (\pi (s ^ {- 1} x) \big)} \\ & {\qquad = \int_ {\mathrm{H}} f (s ^ {- 1} x \xi) d \beta (\xi) = \int_ {\mathrm{H}} \big (\boldsymbol {\gamma} _ {\mathrm{G}} (s) f \big) (x \xi) d \beta (\xi) = \big (\boldsymbol {\gamma} _ {\mathrm{G}} (s) f \big) ^ {b} \big (\pi (x) \big),} \end{array}
\]

whence formula (15), which implies formula (16).

Lemma 4. --- Let \(\lambda\) be a measure \(\neq 0\) on G/H, and \(\mu\) a left Haar measure on G. The following properties are equivalent:

\begin{enumerate}
\def\labelenumi{\alph{enumi})}
\item
  \(\lambda\) is quasi-invariant under G;
\item
  for a subset A of G/H to be locally \(\lambda\) -negligible, it is necessary and sufficient that \(\overline{\pi}^{1}(A)\) be locally \(\mu\) -negligible;
\item
  the measure \(\lambda^{\sharp}\) is equivalent to \(\mu\) .
\end{enumerate}

Suppose this to be the case and let \(\lambda^{\sharp} = \rho \cdot \mu\) , where \(\rho\) is a locally \(\mu\) -integrable function that is everywhere nonzero. Then, for every \(s \in \mathbf{G}\) , the density \(\theta_s\) of \(\gamma_{\mathrm{G / H}}(s)\lambda\) with respect to \(\lambda\) is such that

\[
\theta_ {s} (\pi (x)) = \frac {\rho (s ^ {- 1} x)}{\rho (x)}\tag{17}
\]

locally \(\mu\) -almost everywhere on G.

\begin{enumerate}
\def\labelenumi{\alph{enumi})}
\setcounter{enumi}{2}
\item
  \(\Rightarrow\) b): This follows at once from Prop. 6 a).
\item
  \(\Rightarrow\) a): If property b) holds, then the set of locally \(\lambda\) -negligible subsets of G/H is invariant under G, thus \(\lambda\) is quasi-invariant under G.
\item
  \(\Rightarrow\) c): Assume \(\lambda\) is quasi-invariant under G; for every \(s\in G\) , \(\lambda\) and \(\pmb{\gamma}_{\mathrm{G / H}}(s)\lambda\) are equivalent, therefore \(\lambda^{\sharp}\) and \(\pmb{\gamma}_{\mathrm{G}}(s)\cdot \lambda^{\sharp} = (\pmb{\gamma}_{\mathrm{G / H}}(s)\cdot \lambda)^{\sharp}\) are equivalent (Cor. 1 of Prop. 6); since \(\lambda^{\sharp}\neq 0\) , \(\lambda^{\sharp}\) is equivalent to \(\mu\) (§1, No.~9, Prop. 11).
\end{enumerate}

Moreover, for every \(s \in \mathbf{G}\) ,

\[
\begin{array}{r l} & {(\theta_ {s} \circ \pi) \cdot \lambda^ {\sharp} = (\theta_ {s} \cdot \lambda) ^ {\sharp} = (\boldsymbol {\gamma} _ {\mathrm{G/H}} (s) \lambda) ^ {\sharp} = \boldsymbol {\gamma} _ {\mathrm{G}} (s) \lambda^ {\sharp}} \\ & {\qquad \qquad \qquad \qquad \qquad = (\boldsymbol {\gamma} _ {\mathrm{G}} (s) \rho) \cdot \mu = \frac {\boldsymbol {\gamma} _ {\mathrm{G}} (s) \rho}{\rho} \cdot \lambda^ {\sharp},} \end{array}
\]

whence (17).

The equivalence of a) and b) implies first the uniqueness result already announced, and even a more precise result:

THEOREM 1. --- Let G be a locally compact group, H a closed subgroup of G.

\begin{enumerate}
\def\labelenumi{\alph{enumi})}
\item
  Any two nonzero quasi-invariant measures on G/H are equivalent; the subsets of G/H locally negligible for these measures are those whose inverse image in G is locally negligible for a Haar measure.
\item
  Let \(\lambda, \lambda'\) be two nonzero quasi-invariant measures on G/H. If, for every \(s \in \mathbf{G}\) , the densities of \(\gamma_{\mathrm{G / H}}(s)\lambda\) with respect to \(\lambda\) and of \(\gamma_{\mathrm{G / H}}(s)\lambda'\) with respect to \(\lambda'\) are equal almost everywhere for \(\lambda\) (or \(\lambda'\) ), then \(\lambda\) and \(\lambda'\) are proportional.
\end{enumerate}

The assertion a) follows at once from Lemma 4. Let \(\lambda\) and \(\lambda'\) be two nonzero quasi-invariant measures satisfying the condition of b). Then, for every \(s \in G\) , the densities of \(\gamma_G(s)\lambda^\sharp\) with respect to \(\lambda^\sharp\) and of \(\gamma_G(s)\lambda'^\sharp\) with respect to \(\lambda'^\sharp\) are equal locally \(\mu\) -almost everywhere, therefore (Lemma 3) \(\lambda^\sharp\) and \(\lambda'^\sharp\) are proportional, hence \(\lambda\) and \(\lambda'\) are proportional.

On the other hand, Lemma 4 reduces the search for nonzero quasi-invariant measures on G/H to that for the measures on G equivalent to

Haar measure and of the form \(\lambda^{\sharp}\) . On this subject we have the following lemma:

Lemma 5. --- Let \(\mu\) be a left Haar measure on G, and \(\rho\) a locally \(\mu\) -integrable function. For \(\rho \cdot \mu\) to be of the form \(\lambda^{\sharp}\) , it is necessary and sufficient that, for every \(\xi \in \mathbf{H}\) ,

\[
\rho (x \xi) = \frac {\Delta_ {\mathrm{H}} (\xi)}{\Delta_ {\mathrm{G}} (\xi)} \rho (x)\tag{18}
\]

locally \(\mu\) -almost everywhere on G.

To say that \(\rho \cdot \mu\) is of the form \(\lambda^{\sharp}\) amounts to saying that, for every \(\xi \in \mathbf{H}\) , \(\delta(\xi)(\rho \cdot \mu) = \Delta_{\mathrm{H}}(\xi)\rho \cdot \mu\) (Prop. 4). Now,

\[
\boldsymbol {\delta} (\xi) (\rho \cdot \mu) = (\boldsymbol {\delta} (\xi) \rho) \cdot (\boldsymbol {\delta} (\xi) \mu) = \Delta_ {G} (\xi) (\boldsymbol {\delta} (\xi) \rho) \cdot \mu ,
\]

whence the lemma.

We can now establish the announced existence result, and even a more precise result:

THEOREM 2. --- Let G be a locally compact group, H a closed subgroup of G, μ a left Haar measure on G, and β a left Haar measure on H.

\begin{enumerate}
\def\labelenumi{\alph{enumi})}
\tightlist
\item
  There exist functions \(\rho\) continuous and \(>0\) on \(\mathbf{G}\) , such that
\end{enumerate}

\[
\rho (x \xi) = \frac {\Delta_ {\mathrm{H}} (\xi)}{\Delta_ {\mathrm{G}} (\xi)} \rho (x)
\]

for all \(x \in \mathbf{G}\) and \(\xi \in \mathbf{H}\) .

\begin{enumerate}
\def\labelenumi{\alph{enumi})}
\setcounter{enumi}{1}
\item
  Given such a function \(\rho\) , one can form the measure \(\lambda = (\rho \cdot \mu) / \beta\) on \(\mathbf{G} / \mathbf{H}\) , and \(\lambda\) is a nonzero positive measure quasi-invariant under \(\mathbf{G}\) .
\item
  For \(s, x\) in \(\mathbf{G}\) , \(\rho(sx) / \rho(x)\) depends only on \(s\) and \(\pi(x)\) , hence defines a function \(\chi\) continuous and \(>0\) on \(\mathbf{G} \times (\mathbf{G} / \mathbf{H})\) such that
\end{enumerate}

\[
\chi (s, \pi (x)) = \frac {\rho (s x)}{\rho (x)}.\tag{19}
\]

Then

\[
\boldsymbol {\gamma} _ {\mathrm{G/H}} (s) \lambda = \chi (s ^ {- 1},.) \cdot \lambda \quad \text {   for   all   } s \in \mathrm{G}.\tag{20}
\]

\begin{enumerate}
\def\labelenumi{\alph{enumi})}
\item
  follows from Prop. 7.
\item
  follows from Lemmas 5 and 4.
\item
  follows from (17).
\end{enumerate}

Remarks. --- 1) One deduces from Remark 1 of No.~3 that the nonzero quasi-invariant measures on G/H are none other than the pseudo-images under \(\pi\) of a Haar measure on G.

*2) If G is a Lie group, we shall see later that the function \(\rho\) of Th. 2 can be chosen to be infinitely differentiable.*

Under the conditions of Th. 2, certain results of Nos. 3 and 4 may be specialized as follows (on taking into account Ch. V, §4, Th. 2 and Prop. 2 for passing from properties relative to \(\mu\) to properties relative to \(\rho \cdot \mu\) ):

\begin{enumerate}
\def\labelenumi{\alph{enumi})}
\item
  Let \(f\) be a \(\mu\) -measurable function on \(G\) , with values in a topological space, constant outside a countable union of \(\mu\) -integrable sets; then, the set of \(\dot{x} \in G / H\) such that the function \(\xi \mapsto f(x\xi)\) is not \(\beta\) -measurable is locally \(\lambda\) -negligible.
\item
  Let \(f\) be a \(\mu\) -measurable function \(\geqslant 0\) on \(G\) , zero outside a countable union of \(\mu\) -integrable sets. Then, the function
\end{enumerate}

\[
\dot {x} \mapsto \int_ {\mathrm{H}} ^ {*} f (x \xi) d \beta (\xi)
\]

on G/H is λ-measurable and

\[
\int_ {\mathrm{G}} ^ {*} f (x) \rho (x) d \mu (x) = \int_ {\mathrm{G} / \mathrm{H}} ^ {*} d \lambda (\dot {x}) \int_ {\mathrm{H}} ^ {*} f (x \xi) d \beta (\xi) \quad (\dot {x} = \pi (x)).
\]

\begin{enumerate}
\def\labelenumi{\alph{enumi})}
\setcounter{enumi}{2}
\tightlist
\item
  Let \(\mathbf{f}\) be a \(\rho \cdot \mu\) -integrable function on \(\mathbf{G}\) , with values in a Banach space or in \(\overline{\mathbf{R}}\) . Then, the set of \(\dot{x} \in \mathrm{G} / \mathrm{H}\) such that \(\xi \mapsto \mathbf{f}(x\xi)\) is not \(\beta\) -integrable is \(\lambda\) -negligible; the function \(\dot{x} \mapsto \int_{\mathrm{H}} \mathbf{f}(x\xi) d\beta(\xi)\) is \(\lambda\) -integrable, and
\end{enumerate}

\[
\int_ {\mathrm{G}} \mathbf {f} (x) \rho (x) d \mu (x) = \int_ {\mathrm{G/H}} d \lambda (\dot {x}) \int_ {\mathrm{H}} \mathbf {f} (x \xi) d \beta (\xi).
\]

\begin{enumerate}
\def\labelenumi{\alph{enumi})}
\setcounter{enumi}{3}
\tightlist
\item
  There exists a continuous function \(h \geqslant 0\) on \(G\) , whose support has compact intersection with the saturation \(KH\) of every compact subset \(K\) of \(G\) , and such that \(\int_{H} h(x\xi) d\beta(\xi) = 1\) for every \(x \in G\) . For a function \(g\) on \(G/H\) to be measurable (resp. locally integrable, essentially integrable, integrable) for \(\lambda\) , it is necessary and sufficient that \(h \cdot (g \circ \pi)\) be so for \(\rho \cdot \mu = \lambda^{\sharp}\) ; and, when \(g\) is essentially integrable for \(\lambda\) , one has
\end{enumerate}

\[
\int_ {\mathrm{G/H}} \mathbf {g} (u) d \lambda (u) = \int_ {\mathrm{G}} h (x) \mathbf {g} (\pi (x)) \rho (x) d \mu (x).
\]

